\section{What is sequencing data?}

Modern molecular biology is, at its core, an information-processing
discipline. The hereditary material of every known living organism is
encoded as a long string over the four-letter alphabet
$\Sigma = \{ \mathtt{A}, \mathtt{C}, \mathtt{G}, \mathtt{T} \}$, the four
nucleotides of DNA. \emph{Sequencing} is the process of reading those
strings from a physical sample. The output of a sequencer is, ultimately,
text -- millions or billions of short overlapping fragments called
\emph{reads}, sampled from a much longer underlying genome. The raw
artifact looks deceptively simple:

\begin{lstlisting}[style=DNA]
> read_42 length=120
ACTGTAGCTGATCGTAGCTAGTCACTGTAGCTGTCGTAGCTAGTCCGTAATCGTAGCTAGTCAAACTGTATCGTAGCTAGTCCCACTCGTAGCTACCCCCCCCCCCACTGTAGCTGATC
\end{lstlisting}

The size of these datasets, however, is anything but simple. The cost
per base of DNA sequencing has dropped roughly six orders of magnitude
over the last twenty years -- faster than Moore's law -- and as a
consequence the public Sequence Read Archive at NCBI alone has
accumulated petabytes of raw read data
(Figure~\ref{fig:seq-cost-vs-size}). Most downstream biological
questions -- ``where in the human genome did this read come from?'',
``what variants does this patient carry?'', ``which species are present
in this metagenomic sample?'' -- ultimately reduce to comparing two
strings of nucleotides and quantifying how similar they are. That
comparison is the problem of \emph{sequence alignment}.

\begin{figure}[ht]
  \centering
  \begin{CacheMeCode}{python, custom include command={\input{\robExtAddCachePathAndName{\robExtFinalHash.pgf}}}, set placeholder eval={__LINEWIDTH__}{\lenToCmNoUnit[in]{\linewidth}}}
    import matplotlib.pyplot as plt
    import matplotlib
    matplotlib.use("pgf")
    matplotlib.rcParams.update({
      "font.family": "serif",
      "font.serif": [],
      "text.usetex": True,
    })
    fig, ax = plt.subplots()
    fig.set_size_inches(__LINEWIDTH__, 0.55*__LINEWIDTH__, forward=True)
    years = list(range(2001, 2024))
    cost_per_mb_usd = [
      5292.39, 3898.64, 2986.20, 1135.70, 766.73, 581.92, 397.09,
      102.13, 8.11, 1.36, 0.42, 0.10, 0.06, 0.045, 0.036, 0.024,
      0.014, 0.012, 0.010, 0.008, 0.006, 0.005, 0.005,
    ]
    ax.semilogy(years, cost_per_mb_usd, marker='o', linewidth=2, color='#107AB0',
                label='\\$ per Mbp (NHGRI)')
    moores_law = [cost_per_mb_usd[0] * (0.5 ** ((y - years[0]) / 1.5))
                  for y in years]
    ax.semilogy(years, moores_law, linewidth=1.2, linestyle='--',
                color='#999999', label="Moore's law (1.5x cheaper / 18 mo)")
    ax.set_xlabel('Year', fontsize=8)
    ax.set_ylabel('USD per megabase (log scale)', fontsize=8)
    ax.set_title('Sequencing cost per megabase, 2001--2023', fontsize=9)
    ax.tick_params(axis='both', labelsize=7)
    ax.legend(fontsize=7, frameon=False, loc='upper right')
    ax.grid(True, alpha=0.3, which='both')
    print(get_filename_from_extension(".pgf"))
    plt.savefig(get_filename_from_extension(".pgf"), bbox_inches="tight")
  \end{CacheMeCode}
  \caption{The collapse of sequencing cost over two decades has been
    faster than Moore's law for several technology generations. The
    figures shown are inflation-unadjusted ``per-megabase'' production
    cost numbers from the NHGRI cost survey. The real-world
    consequence is that data volume has grown faster than the ability
    of analysis software to keep up.}
  \label{fig:seq-cost-vs-size}
\end{figure}

\section{Sequence alignment}

Alignment is the operation of placing two strings $s_1, s_2 \in
\Sigma^*$ side-by-side and inserting a special gap character
$\mathtt{-}$ into either string so that, after padding, both have the
same length and no position contains a gap in both strings
simultaneously.  The placement must maximise (or, equivalently,
minimise) some scoring function that rewards matches and penalises
mismatches and gaps.  As a small example, an alignment of
\texttt{GATCAGATA} against \texttt{CATTCAATCA} is

\begin{verbatim}
G A T - C A G A T - A
C A T T C A - A T C A
\end{verbatim}

Aligning two strings of length $n$ and $m$ optimally under an
arbitrary additive scoring function takes $O(nm)$ time and $O(\min(n,
m))$ working memory, using the
Needleman--Wunsch~\cite{needleman_1970}\footnote{Or, depending on the
score model, Smith--Waterman or its 2-piece-affine extension by Suzuki
and Kasahara, used by minimap2.}  dynamic-programming recurrence.
That is acceptable when both strings are short -- a few hundred bases.
It is fatal once we move to whole-genome scale, where a single
$n \approx 3 \cdot 10^9$ human reference sequence aligned against a
$m \approx 10^4$ long-read query already requires
$3 \cdot 10^{13}$ DP-cell evaluations.  That alignment alone would
take hours on a modern CPU.

In practice, exact full-DP alignment is therefore never performed at
genome scale.  Instead, all modern aligners build their pipeline
around a three-stage \emph{seed--chain--extend}
(or \emph{seed-and-extend})
strategy, illustrated in Figure~\ref{fig:seed-chain-extend}.  Most of
the engineering work in this thesis -- and, in particular, all of
Chapter~\ref{r:{minimizers}} -- happens in the seed stage and
the chain stage, where the asymptotically expensive DP work is
\emph{avoided} by spending a smaller amount of time on cheap exact
matching first.

\begin{figure}[ht]
  \centering
  \begin{tikzpicture}[scale=.7, transform shape, node distance=1cm, >=Latex]
    \node(LT)[bigbox]{Load target};
    \node(TI)[bigbox, right=1.5cm of LT]{Build target index};
    \node(LQ)[bigbox, right=1.5cm of TI]{Build query index};
    \node(S)[bigbox, below=of LQ]{Find seeds (exact matches)};
    \node(C)[bigbox, left=of S]{Chain seeds (collinearity test)};
    \node(E)[bigbox, left=of C]{Extend (DP alignment)};
    \node(R)[bigbox, below=of E]{Report alignment};

    \draw[->](LT)--(TI);
    \draw[->](TI)--(LQ);
    \draw[->](LQ)--(S);
    \draw[->](S)--(C);
    \draw[->](C)--(E);
    \draw[->](E)--(R);
  \end{tikzpicture}
  \caption{The seed--chain--extend pipeline shared by minimap2,
    minimap, BWA-MEM, BLASR, Winnowmap, and the bioseq aligner used in
    this work. \emph{Seeds} are short exact-match anchors found via a
    minimizer index. \emph{Chains} are collinear sets of seeds
    consistent with a single underlying alignment. \emph{Extension}
    closes the gaps between adjacent chained seeds with full
    DP-based alignment. The minimizer-collection step (the topic of
    Chapter~\ref{r:{minimizers}}) lives inside ``Build target
    index'' and ``Build query index''.}
  \label{fig:seed-chain-extend}
\end{figure}

\subsection{The seed stage}

A \emph{seed} is a short exact match between query and reference --
typically 15 to 25 bases long. Because we need exact matches, seeds
can be found in $O(N + M)$ time using a hash-table index keyed on
$k$-mer hashes (or, equivalently, an FM-index / suffix array). The
critical observation is that we do not need to find \emph{every}
$k$-mer match: it is sufficient to find $k$-mer matches for a
\emph{representative subset} of the query and reference. This
representative subset is chosen by a sketching scheme called
\emph{minimizers}.

A $(w, k)$\emph{-minimizer} of a sequence is, conceptually, the
lexicographically (or hash-) smallest $k$-mer in every sliding window
of $w$ consecutive $k$-mers. The key property is that the same
$k$-mer is selected from any contiguous substring that contains it as
its window minimum -- so as long as the query and the reference share
a substantial run of identical bases, they are very likely to share a
minimizer position there. In expectation, a $(w, k)$-minimizer scheme
selects roughly $2/(w+1)$ of all $k$-mer positions, dramatically
reducing the index size and the seed-collection cost. With minimap2's
defaults of $k = 15, w = 10$, the index is roughly $1/6$ of the
sequence.

\begin{figure}[ht]
  \centering
  \begin{tikzpicture}[scale=.6, transform shape, node distance=1cm, >=Latex]
    \node (lbl) at (0,0.5) {\textbf{$(w=4, k=5)$-minimizers of GCTATTACCTTAATGTGATGGAC}};

    \node (kbox) at (0,-0.4) {\small Sliding window of $w=4$ consecutive $5$-mers};

    \foreach \i/\letter in {1/G,2/C,3/T,4/A,5/T,6/T,7/A,8/C,9/C,10/T,11/T,12/A,13/A,14/T,15/G,16/T,17/G,18/A,19/T,20/G,21/G,22/A,23/C}
        \node[minimum width=4mm,inner sep=0pt,font=\ttfamily] (b\i) at (\i*4mm-44mm,-1.2) {\letter};

    \draw[rounded corners=1pt, thick, draw=NavyBlue] (b1.south west) rectangle (b8.north east);
    \node[font=\scriptsize, color=NavyBlue, anchor=north] at ($(b4.south)+(0,-0.05)$) {window};

    \draw[rounded corners=1pt, thick, draw=ForestGreen] (b3.south west) rectangle (b7.north east);
    \node[font=\scriptsize, color=ForestGreen, anchor=south] at ($(b5.north)+(0,0.05)$) {minimizer};
  \end{tikzpicture}
  \caption{A toy example of a $(w=4, k=5)$-minimizer. In the window
    spanning positions 1--8 (eight bases = four overlapping 5-mers),
    the chosen minimizer is the 5-mer \texttt{TATTA} starting at
    position 3 -- it has the smallest hash among the four 5-mers
    \texttt{GCTAT, CTATT, TATTA, ATTAC}. The next window (positions
    2--9) shares the same minimizer position (3), so we deduplicate
    and emit it only once.}
  \label{fig:minimizer-toy}
\end{figure}

A more concrete view of how seeds turn into anchor matches is shown
in Figure~\ref{fig:seed-extend-detail}. For each window of $w$
consecutive $k$-mers, the aligner hashes the $k$-mers, picks the
window-minimum, and stores
$(\text{hash},\, \text{position},\, \text{strand})$ in a flat hash
table. Querying then proceeds analogously on the read, and every
match between query and reference is a \emph{seed} that is consistent
with an underlying alignment.

\begin{figure}[ht]
  \centering
  \begin{tikzpicture}[scale=.55, transform shape, node distance=0.5cm, >=Latex]
    \node (qtitle) at (0,4) {\textbf{Query / read}};
    \node[fontscale=-1] (q) at (0,3) {\texttt{...G\hl{CTATTAC}CCATA\hl{TGGAC}TTGCA...}};

    \node (rtitle) at (0,1.5) {\textbf{Reference / target}};
    \node[fontscale=-1] (r) at (0,0.5) {\texttt{...AAGGT\hl{CTATTAC}AAATGCTTGGCAAATTAAA\hl{TGGAC}AT...}};

    \draw[->,>=Latex,color=Plum,thick] (-2.5,2.5) -- (-2.7,1.0);
    \draw[->,>=Latex,color=OrangeRed,thick] (1.5,2.5) -- (3.7,1.0);

    \node[fontscale=-2] at (-2.5, -0.4) {seed 1};
    \node[fontscale=-2] at (3.5,  -0.4) {seed 2};

    \node[align=center, fontscale=-1] at (0,-2) {Two co-linear seeds anchor the alignment.\\Chaining links them; DP closes the gap between them.};

  \end{tikzpicture}
  \caption{From minimizers to anchored alignments. Two minimizer hits
    between query and reference define two \emph{anchors}; chaining
    decides that they are \emph{collinear} (i.e.\ consistent with a
    single underlying alignment), and the extension step closes the
    in-between gap with full dynamic programming. Seed-and-extend
    pipelines exploit the fact that the gap between collinear anchors
    is short, so the DP cost is bounded.}
  \label{fig:seed-extend-detail}
\end{figure}

\subsection{The chaining stage}

A chain is a sequence of anchors $(x_1, y_1), \ldots, (x_n, y_n)$
satisfying $x_1 < x_2 < \ldots < x_n$ and $y_1 < y_2 < \ldots < y_n$
on the reference and query, plus some bound on
$|(x_{i+1} - x_i) - (y_{i+1} - y_i)|$ that captures the maximum
indel size we are willing to tolerate. Chaining is done by a
dynamic-programming recurrence on the anchor list:

\begin{equation*}
f(i) = \max\Big\{\, w_i,\, \max_{j < i}\big[\, f(j) + \alpha(j, i) - \beta(j, i)\,\big] \Big\}
\end{equation*}

where $w_i$ is the anchor weight, $\alpha(j, i)$ rewards co-linearity
and $\beta(j, i)$ penalises gaps. Naively this is $O(N^2)$ in the
number of anchors $N$, but in practice a per-anchor backward look
limited to $h \approx 50$ predecessors gives $O(hN)$ chaining
without measurable loss in chain quality.

The chains are then ranked by score, the best ones are flagged as
\emph{primary}, lower-scoring overlapping chains are flagged as
\emph{secondary}, and the alignment of each primary chain is computed
by closing the per-anchor-pair gaps with DP.

\subsection{The extension stage}

For each adjacent pair of anchors $(x_j, y_j)$ and $(x_{j+1},
y_{j+1})$ in a primary chain, the DP aligner is asked to align the
substring of the reference between $x_j + 1$ and $x_{j+1}$ against
the substring of the query between $y_j + 1$ and $y_{j+1}$. Since the
chain has already constrained both substrings to short lengths
(typically a few hundred bases), the DP cost is bounded.

This is where the bulk of the actual alignment time is spent in
\emph{base-level} aligners; in our bench it accounts for roughly half
of the total wall-clock per pair on long sequences (see
Chapter~\ref{r:{minimizers}}). However, for many downstream
applications (read mapping, taxonomic classification, structural
variant calling on a coarse scale), the chain by itself is enough --
the DP step is skipped or replaced with a divergence
\emph{estimate} computed from the seeded $k$-mer count.

\section{Why this thesis cares about minimizers}

The arithmetic of seed--chain--extend tells us that, among the three
stages, the one most likely to dominate the wall-clock cost depends
strongly on the sequence length:

\begin{itemize}
\item For \textbf{short reads} (~150 bp Illumina-style), the seed
  collection cost is amortised against a deeply-cached reference
  index, the chain has just a handful of seeds, and the DP cost is
  trivial. The aligner is bottlenecked on memory-access latency.
\item For \textbf{long reads} (PacBio HiFi / Oxford Nanopore,
  $10^4$--$10^6$ bp), the chain becomes large -- tens of thousands of
  anchors per pair -- and the DP-extension step dominates.
\item For \textbf{very-long-read overlap mode}
  (PacBio HiFi, ${\sim}10^6$ bp pairs aligned to references of similar
  scale), \emph{everything} matters. The seed step has to chew
  through millions of $k$-mer hashes, the chain has 100~k+ anchors,
  and the alignment is cumulatively a few hundred milliseconds per
  pair.
\end{itemize}

For the very-long-read regime, even a 25\,\% reduction in the time
spent computing minimizers translates into a perceptible improvement
in throughput. Worse, the minimizer step is the only stage in the
pipeline whose runtime cost is structurally \emph{linear} in the
input length and not dominated by alignment depth -- it sees every
single base. Anything that touches every base benefits massively from
SIMD vectorisation. This thesis devotes
Chapter~\ref{r:{minimizers}} to making the SIMD minimizer
collector competitive with -- and ultimately faster than -- the
hand-tuned scalar implementation in minimap2.

\section{The bioseq pipeline}

The bioseq aligner used throughout this thesis is implemented as a
PostgreSQL extension. The on-disk layout, the index build, and the
SQL surface are described in
Chapter~\ref{r:{postgres}}; from the algorithmic point of
view, however, the bioseq aligner is a faithful port of the minimap2
seed--chain--extend pipeline, with three deliberately selectable
implementations of each stage:

\begin{description}
\item[Minimizer] \textsc{scalar} (an in-house scalar
  port of minimap2's \texttt{mm\_sketch}); \textsc{simd} (a literal
  SIMDification of the rescan loop); \textsc{two-stack v8/v9/v11/v12/v13/v14}
  (the family of SIMD-Minimizers-inspired variants developed in
  Chapter~\ref{r:{minimizers}}).
\item[Chaining] \textsc{scalar-mm2} (the minimap2 backwards-DP
  $h{=}50$ chainer ported to C++); \textsc{rmq} (range-minimum-query
  chainer, $O(N \log N)$); \textsc{simd-mm2} (a 4-lane batched DP
  chainer); \textsc{simd-mm2-v2} (a dual-batch unrolled chainer
  inspired by the same loop-restructuring technique that worked for
  minimizers, but which turns out not to transfer to chaining --
  see~\S\ref{sec:v15-chain-v2}).
\item[Base aligner] \textsc{ksw2} (the same Suzuki--Kasahara DP
  used by minimap2); \textsc{block} (a block-based aligner inspired
  by Block-Aligner~\cite{block_aligner_paper}).
\end{description}

The combination grid yields up to $3 \times 3 \times 2 = 18$
algorithm configurations per benchmark run -- the actual bench
filters this down to the dozen-odd configurations that are worth
comparing at any given length bucket. The benchmark harness emits a
machine-readable \texttt{report.json} that the
\path{bioseq/RESEARCH/BENCHMARKING_SPEEDUP/view\_report.ipynb} notebook
loads to produce all the comparison charts surfaced in the
engineering reports. The contents of that notebook -- the pre-tier-0
overview, the seven-tier headline-to-deep-dive structure, and the
new minimizer-variant comparison section added for this thesis --
mirror the structure of Chapter~\ref{r:{minimizers}}.

\section{Roadmap of the rest of the thesis}

\begin{itemize}
\item Chapter~\ref{r:{minimizers}} is the engineering core of
  this thesis. It walks through the iterative work to produce a
  SIMD-friendly minimizer collector that beats scalar minimap2 by
  $\geq 10\,\%$ on the user-stated very-long-read regime
  ($1$--$5$\,Mbp), starting from the original minimap2 algorithm,
  walking through the naive SIMDification (which lost), porting in
  the two-stacks design from the SIMD-Minimizers paper of Groot
  Koerkamp \& Martayan~\cite{groot_koerkamp_simd_minimizers_2025},
  and applying loop unrolling, dual-stream interleaving, and forced
  inlining to bridge the remaining gap. It also documents the
  surprising downstream effects on \texttt{mm\_est\_err} (the
  divergence estimator) and on the seed-hit collection step --
  fixing those was load-bearing for hitting the
  $\geq 10\,\%$ goal.
\item Chapter~\ref{r:{postgres}} covers the PostgreSQL
  integration: how genomic sequences are stored as TOASTed columns,
  how the minimizer index is exposed as a SQL function, and how
  queries plan against it.
\end{itemize}
